% ACCESS supplement request, PHY250179: progress report and plan.
% Upload this file (and the figures named in \figpanel) to Overleaf.
% Bracketed items in bold are for the PI to fill in. Figures: put the
% files next to this .tex; a missing file shows a placeholder box, so the
% document compiles before the figures are added.
\documentclass[11pt]{article}
\usepackage[letterpaper,margin=0.85in]{geometry}
\usepackage{graphicx}
\usepackage{booktabs}
\usepackage{amsmath,amssymb}
\usepackage{xcolor}
\usepackage[hidelinks]{hyperref}
\usepackage{titlesec}
\usepackage{enumitem}
\titlespacing*{\section}{0pt}{1.0ex}{0.4ex}
\titleformat{\section}{\normalfont\large\bfseries}{\thesection.}{0.5em}{}
\setlist{nosep,leftmargin=1.4em}
\setlength{\parskip}{0.35ex}
\setlength{\parindent}{1.2em}
\pagestyle{empty}

\newcommand{\todo}[1]{\textbf{[#1]}}
% \figpanel{file}{caption}: a panel of a 2 x 2 figure block; a placeholder
% box if the file is missing.
\newcommand{\figpanel}[2]{%
  \begin{minipage}[t]{0.49\textwidth}\centering
    \IfFileExists{#1}{\includegraphics[width=\linewidth,height=5.2cm,keepaspectratio]{#1}}%
      {\fbox{\parbox[c][5.0cm][c]{0.95\linewidth}{\centering\small\texttt{\detokenize{#1}}\\(figure to be added)}}}\par\vspace{0.3ex}
    {\footnotesize #2}
  \end{minipage}}

\begin{document}

\begin{center}
{\large\bfseries Supplement request for PHY250179:\\
High-resolution cosmological simulations of intergalactic gas}\\[0.4ex]
Matthew McQuinn (PI), University of Washington \quad $\cdot$ \quad NASA grant 80NSSC24K1220 \quad $\cdot$ \quad \todo{date}\\
Accelerate project; supplement of 1{,}500{,}000 ACCESS credits requested.
Resources used so far: \todo{SDSC Expanse; other systems}; credits spent / remaining: \todo{numbers}.
\end{center}

\section{Progress to date}

The allocation supported high-resolution simulations of the intergalactic
medium (IGM) at $z = 3$--$7$ in two directions: the response of IGM gas to
reionization at the kiloparsec scale, and the structure of cold gas in the
cosmic web. Four papers resulted, three on arXiv and one about to be
submitted (Fig.~\ref{fig:results}).

\textbf{(a) SAGUARO, a suite of IGM reionization simulations} (Cain, Das,
D'Aloisio, Foreman, Scannapieco, Moreno, Lugatiman, Cohon, Tohfa \& Trac 2026,
arXiv:2603.25788). Over two hundred coupled radiative-hydrodynamics
simulations of IGM gas during and after reionization in boxes of 2 and
0.25~$h^{-1}$Mpc, at resolutions that resolve the Jeans scale of the cold
neutral IGM, spanning a grid of photoionization rates, reionization redshifts
and box-scale densities, with X-ray preheating, recombination radiation,
baryon--dark-matter streaming and alternative dark-matter cosmologies. The
paper sets out the suite and its first results on the gas dynamics and
thermal structure, the ionizing opacity and self-shielding, the IGM's drain on
the reionization photon budget, the halo mass function and Ly$\alpha$
transmission. \todo{which runs used this allocation, and their credits}

\textbf{(b) An emulator for the ionizing-photon mean free path} (Tohfa, Cain,
McQuinn \& D'Aloisio 2026, arXiv:2602.03923; JCAP 06 (2026) 094). A
deep-learning emulator trained on high-resolution radiative-transfer
simulations of ionization fronts in $\approx 2\,h^{-1}$Mpc volumes, which
capture the kiloparsec-scale clumping that large reionization volumes cannot
resolve. Combined with the measured photoionization rate and mean free path at
$z = 5$--$6$, and without power-law assumptions, it implies a factor 2--3
decline in the ionizing emissivity between $z = 6$ and $4.8$; it also serves
as a sub-grid opacity for large simulations.

\textbf{(c) A self-consistent analytic model of the photoionization rate and
the reionization history} (Cain, Croker, D'Aloisio, Georgiev, Tohfa, Zhu \&
Windhorst 2026, arXiv:2606.19449). An analytic formalism derived from the
cosmological radiative-transfer equation that predicts the ionized fraction
and $\Gamma_{\rm HI}$ together, tested against the radiative-transfer
simulations: percent-level in the ionized fraction and 20--30\,\% in
$\Gamma_{\rm HI}$ at $z \lesssim 6$, which makes the forest's
$\Gamma_{\rm HI}$ usable in fast inference of the reionization history.

\textbf{(d) Multiphase gas in high-$z$ cosmic sheets} (Willard, Faerman,
McCarty, McQuinn, Mandelker \& Quinn, to be submitted). A suite of moving-mesh
(AREPO) simulations of individual $z = 3$--$6$ cosmic sheets, 4 and 8~cMpc/$h$
boxes at $2\times128^3$ to $2\times1024^3$ particles (cold-gas cells of
0.4--3~pkpc), with an SPH comparison suite in MP-Gadget; all runs at the San
Diego Supercomputer Center. The sheets' cold cores fragment through the
nonlinear thin-shell instability and then Kelvin--Helmholtz, not the thermal
instability proposed before; fragmentation grows with resolution and is not
converged even at the highest; the covering fractions of partial Lyman-limit
and Lyman-limit systems rise by factors of 2--4 across the resolution range;
and cosmic sheets plausibly supply of order 1\,\% of LLSs and 10\,\% of pLLSs
at $z = 4$, a large share of the metal-free absorbers. \todo{credits used}

\begin{figure}[t]
\centering
\figpanel{fig_saguaro.pdf}{(a) SAGUARO: \todo{what the panel shows}.}\hfill
\figpanel{fig_emulator.pdf}{(b) The mean-free-path emulator: \todo{what the panel shows}.}\\[0.8ex]
\figpanel{fig_analytic.pdf}{(c) The analytic $\Gamma_{\rm HI}$ and $x_{\rm HII}$ model: \todo{what the panel shows}.}\hfill
\figpanel{fig_sheets.pdf}{(d) Cosmic sheets: $N_{\rm HI}$ of the high-$T_{\rm vir}$ sheet across resolutions (Willard et al.).}
\caption{One result from each paper the allocation produced.}
\label{fig:results}
\end{figure}

\textbf{Software built for the next phase.} A C++ port of Trac's RadHydro (the
code behind (a)--(c)), validated bit for bit against the original, now runs
with MPI\,+\,OpenMP through AMReX on 192-core nodes ($512^3$ reionization runs
on five nodes, deterministic across rank counts), with adaptive mesh refinement
for gas, radiation and gravity validated to two levels, a star-formation sink,
and a first refined $512^3$ production run whose cost is measured. This
development ran on the group's own cluster; the supplement funds the runs it
enables.

\section{Why more credits: turbulence driven by real ionization fronts}

Cain, McQuinn, Scannapieco, D'Aloisio \& Trac (2026, PRD 113, 023050) showed
with uniform 0.25~$h^{-1}$Mpc boxes at $1024^3$ and $2048^3$ that the heating
by ionization fronts drives kiloparsec-scale turbulence once cells are smaller
than $\approx 0.5$~kpc/$h$, and that this turbulence can grow volume-filling
magnetic fields to the strengths the TeV-blazar limits require. Those boxes
were ionized by plane-parallel fronts at a chosen redshift. The new code makes
two things possible that were not when the project was written:
\begin{enumerate}
\item \textbf{The same boxes in the new code (C2):} $1024^3$ at 244~pc/$h$ with
  plane-parallel fronts, then preheating ($T_{\min} = 100$ and 1000~K) and
  reionization-redshift variants: the code-to-code check and the baseline.
\item \textbf{Zoom-in runs of one region of a full reionization simulation
  (Z2, Z3):} a 40~$h^{-1}$Mpc, $512^3$ box with sources and escape fraction
  calibrated to the Ly$\alpha$ forest, in which one region of
  1.25~$h^{-1}$Mpc ($1024^3$ cells at 1.2~kpc/$h$) and then 0.31~$h^{-1}$Mpc
  ($1024^3$ at 0.3~kpc/$h$) is refined by six and eight levels, with initial
  conditions seeded at the region's resolution. The region then meets the real
  front: its arrival time, speed, direction and spectrum, and the surrounding
  filaments and bulk flow. Regions are chosen by density and by ionization
  time. This gives the turbulence driven by a real front and the first
  self-consistent setting for the Biermann battery at the front.
\end{enumerate}

\section{Resources requested}

\todo{Resource and the conversion from the ACCESS exchange calculator.} From
the measured costs of this code on 192-core AMD Genoa nodes (one hydro
cell-update $\approx 3\,\mu$s with refinement, one ray-substep
$\approx 2\,\mu$s) and the refined $512^3$ production run ($2.6\times$ the
unrefined step time):

\begin{center}\small
\begin{tabular}{lll}
\toprule
Run & Finest-level cells & Core-hours \\
\midrule
C2: $1024^3$, 0.25~$h^{-1}$Mpc, $z = 300 \to 4$ & $1.1\times10^9$ & $\approx 1\times10^5$ \\
C2 variants: $T_{\min} = 100$, 1000~K; a second $z_{\rm re}$ & & $\approx 3\times10^5$ \\
Z2: $512^3$ base + 6 levels, $1024^3$ over 1.25~$h^{-1}$Mpc & $1.1\times10^9$ & $\approx 3\times10^4$ \\
Z3: the same region, 8 levels, $1024^3$ over 0.31~$h^{-1}$Mpc & $1.1\times10^9$ & $\approx 1\times10^5$ \\
A second Z3 region (a different ionization time) & $1.1\times10^9$ & $\approx 1\times10^5$ \\
Resolution ladders ($64^3$--$512^3$) and analysis & & $\approx 1\times10^5$ \\
\midrule
Total & & $\approx 7\times10^5$ core-hours \\
\bottomrule
\end{tabular}
\end{center}
\noindent Requested as 1.5 million credits with the conversion and a margin
for the refinement overhead. Memory under 1~TB per run; checkpoints
$\approx 0.4$~TB every ten steps; the finest-level cubes 4~GB per field per
snapshot. Storage: \todo{amount and system}.

\section{Timeline}

\todo{Months 1--2: C2 and variants. Months 2--5: Z2, Z3, the second region.
Month 6: analysis, and the renewal request for the full suite (several regions
$\times$ preheating, the MHD runs).}

\end{document}
